Foram utilizados três conjuntos de dados: Energy Efficiency (regressão), Wine Quality (Red) e Iris (classificação multiclasse). O pré-processamento incluiu normalização dos atributos, codificação de variáveis categóricas e divisão dos dados em treino e teste. A escolha da normalização se justifica pela diferença de escala entre os atributos, impactando algoritmos baseados em distância e gradiente.

Os algoritmos avaliados foram: Regressão Linear, Random Forest, SVR e MLP para regressão; Regressão Logística, Random Forest, SVC e MLP para classificação. Cada setup foi executado com 30 repetições para algoritmos estocásticos, variando a semente de inicialização. As métricas utilizadas foram RMSE, MAE, R² para regressão e acurácia, F1-score, Kappa para classificação.
